\chapter{Economics of Informat ion}


\section{Adverse Selection (Hidden Knowledge)}

Information is closely related to uncertainty.

Information determines the degree of uncertainty that an individual has about an event.

An information set represents the specific information held by an individual.

Individuals facing uncertainty want to incorporate all available information into their decision-making process.

This means that an individual's action is generally a function of his or her information set.

Allowing actions to depend on information creates substantial complexity, because different individuals may possess different information structures.

Suppose individual \(A\)'s action depends on his information set
\[
I_A.
\]

Other agents may then infer something about
\[
I_A
\]
from observing \(A\)'s behavior.

Therefore, the information available to other agents depends, through observed actions, on \(A\)'s information set.

As a result, individual \(B\)'s information structure is influenced by \(A\)'s actions, and similarly \(A\)'s information structure may be influenced by \(B\)'s actions.

This type of interaction generally requires the tools of game theory.

If all agents possess equally limited information, then no one has an informational advantage over others.

In contrast, asymmetric information generates problems of opportunistic behavior, in which the informed person benefits at the expense of the less informed person.

The informed individual is often called the \emph{agent}, while the less informed individual is called the \emph{principal}.

The two major forms of opportunistic behavior are:

\begin{enumerate}
    \item adverse selection,
    \item moral hazard.
\end{enumerate}

Adverse selection refers to opportunistic behavior arising from hidden knowledge.

More precisely, adverse selection occurs when an informed person benefits from trading or contracting with a less informed person who cannot observe some important characteristic of the informed person.

Thus, adverse selection is fundamentally a problem of hidden information before a contract is made.

A classic example is the market for life insurance.

Individuals purchasing life insurance policies generally possess better information about their own health conditions than insurance companies do.

If an insurance company offers insurance at a uniform premium rate, individuals with higher mortality risk are more likely to purchase insurance.

As a result, the pool of insured individuals becomes disproportionately composed of high-risk consumers.

This raises the average expected payout faced by the insurance company.

If the insurer cannot distinguish high-risk individuals from low-risk individuals, the premium must rise to cover expected losses.

However, as premiums rise, low-risk individuals may choose not to purchase insurance because the premium exceeds their expected benefit.

Consequently, the composition of insured individuals deteriorates further, leaving an even riskier insurance pool.

This process is sometimes called the \emph{adverse selection spiral}.

In extreme cases, the market itself may collapse because only the highest-risk individuals remain willing to purchase insurance.

\subsection{The Market for Lemons}

One of the most famous examples of adverse selection is Akerlof's ``Market for Lemons.''

Suppose there are two types of used cars:

\begin{itemize}
    \item high-quality cars,
    \item low-quality cars (lemons).
\end{itemize}

Sellers know the quality of their own cars, but buyers cannot perfectly distinguish quality before purchase.

Because buyers face uncertainty, they are willing to pay only an average price reflecting expected quality.

However, owners of high-quality cars may refuse to sell at this average price because it is too low relative to the true value of their cars.

As high-quality sellers leave the market, the average quality of cars offered for sale declines.

Rational buyers anticipate this decline and lower the price they are willing to pay.

This further drives high-quality sellers out of the market.

Eventually, only low-quality cars remain in the market.

Thus, asymmetric information can reduce market efficiency and may even eliminate mutually beneficial trades.

\subsection{Signaling and Screening}

Markets often develop mechanisms to reduce adverse selection.

Two important mechanisms are:

\begin{enumerate}
    \item signaling,
    \item screening.
\end{enumerate}

\paragraph{Signaling}

Signaling occurs when informed agents take actions that credibly reveal their private information.

For example, educational attainment may serve as a signal of worker productivity.

If obtaining education is less costly for high-productivity workers than for low-productivity workers, then education can separate different worker types.

Employers may then use educational credentials as indicators of productivity.

\paragraph{Screening}

Screening occurs when the uninformed side designs contracts or mechanisms to induce self-selection.

For example, insurance companies may offer multiple insurance contracts:

\begin{itemize}
    \item a high-premium, low-deductible contract,
    \item a low-premium, high-deductible contract.
\end{itemize}

High-risk individuals tend to choose more comprehensive coverage, while low-risk individuals choose cheaper contracts with higher deductibles.

Thus, different types reveal themselves through their choices.

\subsection{Efficiency Implications}

Adverse selection creates inefficiency because mutually beneficial trades may fail to occur.

The key problem is that prices no longer reflect only the value of goods or services; they also reflect informational asymmetries.

Asymmetric information may therefore lead to:

\begin{itemize}
    \item market breakdown,
    \item inefficient allocation,
    \item excessive precautionary pricing,
    \item reduced trade volume,
    \item welfare losses.
\end{itemize}

Much of modern information economics studies how institutions, contracts, reputation systems, regulation, and market design can reduce these inefficiencies.




If an insurance company offers insurance at a premium based on the average risk of the population, then a disproportionately large share of unhealthy individuals will purchase the policy.

Because of this adverse selection, the insurance company will pay claims on more policies than it would if healthy and unhealthy individuals purchased insurance in proportion to their shares in the overall population.

Thus, adverse selection creates a market failure by reducing the size of the market or potentially eliminating it altogether, thereby preventing mutually beneficial transactions.

\bigskip

Moral hazard refers to opportunistic behavior arising from hidden actions.

More precisely, moral hazard occurs when an informed person takes actions that cannot be perfectly observed by another party.

For example, an employee may shirk and exert insufficient effort if the employer cannot monitor performance perfectly.

Similarly, insured individuals may engage in risky behavior after obtaining insurance coverage because part of the cost of risky actions is borne by the insurance company.

They may also fail to take appropriate precautions that would reduce the probability of accidents or losses.

Examples include:

\begin{itemize}
    \item reckless driving after obtaining automobile insurance,
    \item reduced care in preventing theft after purchasing theft insurance,
    \item reduced workplace effort when monitoring is weak.
\end{itemize}

Moral hazard leads to inefficiency because hidden actions reduce output or increase the frequency of accidents and losses.

Therefore, moral hazard is another important source of market failure.

\section{Insurance Market}

\subsection{Rothschild--Stiglitz Model}

There exists an average health risk in the population, but risk differs across individuals.

Suppose individuals possess better information about their own risk type than insurance companies do.

What happens if an insurance company offers insurance at a premium designed to break even at the average population risk?

Individuals in low-risk groups may decide not to purchase insurance because the premium is too high relative to their expected losses.

In contrast, individuals in high-risk groups are more likely to purchase insurance because the premium is attractive relative to their expected losses.

Consequently, the average risk of the insured pool becomes higher than the average risk of the overall population.

As a result, the insurance company experiences losses.

This phenomenon is adverse selection.

More generally, adverse selection describes a situation in which agents possess different information before participating in the market.

\subsection{Equilibrium in Insurance Markets}

An equilibrium in an insurance market must satisfy the condition that no insurance company can profitably introduce a new contract that attracts consumers away from the equilibrium contracts.

In other words, there must not exist an alternative contract that:

\begin{itemize}
    \item yields positive expected profits,
    \item and is preferred by some consumers to the contracts already offered.
\end{itemize}

\subsection{Equilibrium with Identical Individuals}

Suppose all individuals have:

\begin{itemize}
    \item identical initial wealth,
    \item identical preferences over wealth,
    \item identical risk probabilities.
\end{itemize}

Let initial wealth be
\[
W.
\]

Each individual either:

\begin{itemize}
    \item suffers no loss,
    \item or suffers a loss of size
    \[
    x.
    \]
\end{itemize}

Without insurance, there are therefore two possible wealth levels:
\[
W_1=W,
\]
and
\[
W_2=W-x.
\]

Now consider an insurance contract characterized by:

\begin{itemize}
    \item insurance coverage
    \[
    y,
    \]
    \item premium
    \[
    h.
    \]
\end{itemize}

If the loss occurs, the insurance company pays
\[
y.
\]

The premium
\[
h
\]
must be paid regardless of whether the loss occurs.

Hence, with insurance, wealth in the two states becomes:
\[
W_1=W-h,
\]
if no loss occurs,
and
\[
W_2=W-x+y-h,
\]
if the loss occurs.

Suppose the probability of loss is
\[
\pi.
\]

Expected utility is therefore
\[
EU
=
(1-\pi)u(W-h)
+
\pi u(W-x+y-h).
\]

If insurance companies are competitive and risk neutral, zero expected profit implies
\[
h=\pi y.
\]

Substituting this condition into expected utility yields
\[
EU
=
(1-\pi)u(W-\pi y)
+
\pi u(W-x+y-\pi y).
\]

The individual chooses insurance coverage
\[
y
\]
to maximize expected utility.

The first-order condition is
\[
-(1-\pi)\pi u'(W-\pi y)
+
\pi(1-\pi)u'(W-x+y-\pi y)
=
0.
\]

Equivalently,
\[
u'(W-\pi y)
=
u'(W-x+y-\pi y).
\]

If utility is strictly concave, then marginal utility is strictly decreasing.

Therefore,
\[
W-\pi y
=
W-x+y-\pi y,
\]
which implies
\[
y=x.
\]

Thus, under actuarially fair insurance and risk aversion, individuals choose full insurance.

With full insurance,
\[
W_1=W_2=W-\pi x.
\]

Insurance completely eliminates risk and equalizes wealth across states.



With the insurance contract, wealth is
\[
W_1=W-h
\]
in the no-loss state, and
\[
W_2=W-h-x+y
\]
in the loss state.

Assume that all individuals face the same probability of loss,
\[
p.
\]

In this case, the profit of the insurance company is
\[
\pi
=
p(h-y)+(1-p)h.
\]

Equivalently,
\[
\pi
=
h-py.
\]

In a perfectly competitive market with constant returns to scale technology, profit must be zero:
\[
\pi=0.
\]

Therefore,
\[
h=py.
\]

In Figure 1, an individual is at point \(A\) without insurance.

The budget constraint is
\[
W-px
=
W^0
=
(1-p)W_1+pW_2.
\]

The fair-odds line is
\[
\frac{dW_2}{dW_1}
=
-\frac{1-p}{p}.
\]

This is the line depicted in Figure 1.

Below this line, the insurance company obtains positive profit.

Purchasing an insurance contract corresponds to moving along the fair-odds line.

The individual's expected utility is
\[
u^0
=
(1-p)u(W_1)+pu(W_2).
\]

Thus, the slope of the indifference curve is
\[
\frac{dW_2}{dW_1}
=
-
\frac{(1-p)u'(W_1)}{pu'(W_2)}.
\]

At the tangency point \(E\), we obtain
\[
\frac{u'(W_1)}{u'(W_2)}=1.
\]

Hence,
\[
W_1=W_2.
\]

This corresponds to full insurance coverage.

\subsection{Equilibrium with Heterogeneous Individuals and Complete Information}

Now assume that there are two types of individuals.

A fraction
\[
\lambda
\]
have a relatively high probability of loss:
\[
p_H,
\]
and a fraction
\[
1-\lambda
\]
have a relatively low probability of loss:
\[
p_L.
\]

Assume
\[
p_L<p_H.
\]

It is also assumed that the principal, namely the insurance company, observes each individual's type.




Since the insurance company can identify individual types, it offers two different fair-odds lines, one for each type.

Type \(H\) chooses point \(B\), while type \(L\) chooses point \(C\).

Thus, an efficient separating equilibrium is obtained.

Both type \(H\) and type \(L\) choose full coverage.

\subsection{Equilibrium with Heterogeneous Individuals and Asymmetric Information}

Now assume that individuals are heterogeneous, but the insurance company cannot observe each individual's type.

This is the case of asymmetric, or incomplete, information.

We consider two types of equilibria:

\begin{enumerate}
    \item a pooling equilibrium, in which both types buy the same contract;
    \item a separating equilibrium, in which the two types buy different contracts.
\end{enumerate}

We first examine the feasibility of a pooling equilibrium.

If everyone purchases the same insurance contract, then the aggregate probability of loss from the insurance company's perspective is
\[
p
=
\lambda p_H+(1-\lambda)p_L.
\]

Hence, the zero-profit condition is
\[
h=py.
\]

Equivalently, the contract must lie on the fair-odds line with slope
\[
-\frac{1-p}{p}.
\]

We denote by \(D\) an arbitrary point on the fair-odds line and test whether it can be an equilibrium.

Construct the indifference curves of type \(H\) and type \(L\) that pass through point \(D\).

At point \(D\), the indifference curve of type \(L\) is steeper than that of type \(H\).

This property is called the \textbf{single-crossing condition}.




The single-crossing condition is
\[
\left|
\frac{dW_2}{dW_1}
\right|_{L}
=
\left|
-
\frac{(1-p_L)u'(W_1)}{p_Lu'(W_2)}
\right|
>
\left|
-
\frac{(1-p_H)u'(W_1)}{p_Hu'(W_2)}
\right|
=
\left|
\frac{dW_2}{dW_1}
\right|_{H}.
\]

\begin{theorem}
There cannot exist a pooling equilibrium.
\end{theorem}

\begin{proof}
Suppose that point \(D\) is a pooling equilibrium.

Such a pooling equilibrium cannot be sustained because type \(L\) individuals would prefer to move to another contract such as point \(B\), or even point \(A\).

Therefore, no pooling equilibrium can survive competition.
\end{proof}

We now examine the feasibility of a separating equilibrium consisting of two different contracts, each purchased by a different group.

In this case, the zero-profit condition can still be maintained if the two contracts lie on their respective fair-odds lines through the initial endowment point \(A\).

Suppose contracts corresponding to points \(B\) and \(C\) are offered.

Then both types would prefer contract \(C\), generating losses for insurance companies.

Hence, the contract intended for type \(L\) must be moved downward along the low-risk fair-odds line until type \(H\) individuals no longer prefer switching away from contract \(B\).

The appropriate point is denoted by \(D\), where the indifference curve of type \(H\) intersects the fair-odds line of type \(L\).

If contracts corresponding to points \(B\) and \(D\) are offered, then:

\begin{itemize}
    \item type \(H\) chooses \(B\),
    \item type \(L\) chooses \(D\).
\end{itemize}

Under this allocation:

\begin{itemize}
    \item both contracts yield zero profit,
    \item neither type has an incentive to switch contracts.
\end{itemize}

Therefore, a separating equilibrium exists.




on these fair-odds lines are separating.

Hence, we derive the following theorem.

\paragraph{Theorem 2.}

Assume that a separating equilibrium exists.

Then the high-risk group receives full insurance, while the low-risk group receives partial insurance.

\paragraph{Exercise 1.}

Let
\[
W_1=100
\]
denote wealth without loss, and let
\[
W_2=60
\]
denote wealth with loss.

\begin{enumerate}
    \item[(A)] Suppose the probability of each state is \(0.5\) for each consumer. Derive the market equilibrium, premium, and coverage.

    \item[(B)] Now suppose there are two types of consumers with
    \[
    p_H=0.8
    \]
    and
    \[
    p_L=0.2.
    \]
    Assume that the insurance company can identify each consumer's type. Derive the market equilibria, premiums, and coverages.

    \item[(C)] Suppose the number of type \(H\) consumers is the same as the number of type \(L\) consumers, but the insurance company cannot identify consumer types.

    The utility function is
    \[
    u(W_i)=\sqrt{W_i},
    \qquad
    i=1,2.
    \]

    Derive the market equilibria, premiums, and coverages.
\end{enumerate}

\subsection{Signaling}

If the principal has only a single choice variable, it is impossible to obtain a separating equilibrium.

An additional variable is needed so that two states in an information set can be distinguished.

An insurance policy has two natural variables:

\begin{itemize}
    \item premium,
    \item coverage level.
\end{itemize}

Therefore, a separating equilibrium may be possible in an insurance market.



equilibrium can be constructed.

When there is originally only a single variable, a signal can serve as the required second variable.

A signal takes the form of an action that has economic consequences.

Observation of the action taken by an agent may reveal information that is otherwise hidden.

The difficulties associated with adverse selection can be overcome only if the costs of signaling differ across types of agents.

\paragraph{Example 1. Job Market Signaling}

Spence (1973) considers a labor market composed of two types of individuals:

\begin{itemize}
    \item type \(L\): low productivity,
    \item type \(H\): high productivity.
\end{itemize}

Let the fraction of type \(L\) workers be
\[
q,
\]
and the fraction of type \(H\) workers be
\[
1-q.
\]

Productivity is measured by marginal product.

The marginal product of type \(L\) is
\[
1,
\]
while the marginal product of type \(H\) is
\[
2.
\]

The average marginal product is therefore
\[
q\cdot 1+(1-q)\cdot 2
=
2-q.
\]

If employers cannot distinguish between worker types, then the equilibrium wage paid to both types is
\[
2-q.
\]

Now suppose that workers can obtain education before entering the labor market.

Let \(y\) denote the amount of education.

The cost of education is
\[
c_L(y)=y
\]
for type \(L\), while
\[
c_H(y)=\frac{1}{2}y
\]
for type \(H\).

Education level is the signal.

Suppose employers choose a critical education level
\[
y^*.
\]

They pay wage
\[
1
\]
to workers with education level below \(y^*\), and wage
\[
2
\]
to workers with education level at or above \(y^*\).

Given this wage schedule, individuals choose \(y\) to maximize net return:
\[
\text{wage}-\text{education cost}.
\]

Assume
\[
1\le y^*\le 2.
\]

Then type \(L\) chooses
\[
y=0,
\]
while type \(H\) chooses
\[
y=y^*.
\]

In equilibrium, firms satisfy the zero-profit condition because each worker is paid his marginal product.

The separating equilibrium is not unique.

Any value of \(y^*\) satisfying
\[
1\le y^*\le 2
\]
yields a separating equilibrium.

\section{A Formal Model of Adverse Selection}

In this model, a principal wants to delegate a task to an agent who has access to information not available to the principal.

To achieve an efficient use of economic resources, the contract between the principal and the agent must elicit the agent's private information.

This is possible only by giving up some information rent to the agent.

\subsection{The Basic Model}

A principal delegates production,
\[
q,
\]
to an agent.

Let
\[
S(q)
\]
be the value of production to the principal.

Assume
\[
S'(q)>0,
\qquad
S''(q)<0,
\qquad
S(0)=0.
\]

The production cost of the agent will depend on the agent's private type.




The production cost of the agent is not observable to the principal.

Let
\[
\Theta=\{\underline{\theta},\overline{\theta}\}
\]
be the set of possible marginal costs, where
\[
\Delta\theta
=
\overline{\theta}-\underline{\theta}
\ge 0.
\]

The agent can be either:

\begin{itemize}
    \item efficient, with marginal cost
    \[
    \underline{\theta},
    \]
    \item inefficient, with marginal cost
    \[
    \overline{\theta}.
    \]
\end{itemize}

The respective probabilities are

\[
\nu
\]
for the efficient type and
\[
1-\nu
\]
for the inefficient type.

Hence, the production cost functions are
\[
c(q,\underline{\theta})
=
\underline{\theta}q
\]
with probability
\[
\nu,
\]
and
\[
c(q,\overline{\theta})
=
\overline{\theta}q
\]
with probability
\[
1-\nu.
\]

The set
\[
\mathcal{A}
=
\{
(q,t)
:
q\in\mathbb{R}_+,
\;
t\in\mathbb{R}
\}
\]
is the set of feasible allocations, where

\begin{itemize}
    \item \(q\) is production,
    \item \(t\) is the transfer paid to the agent.
\end{itemize}

These variables are observable and verifiable by a third party.

\subsection{The Complete Information Optimal Contract}

Under complete information, efficient production levels are obtained by equating the principal's marginal value and the agent's marginal cost:
\[
S'(q^*_L)=\underline{\theta},
\]
and
\[
S'(q^*_H)=\overline{\theta}.
\]

We assume that production is socially valuable.

Thus,
\[
W^*
=
S(q^*_H)-\overline{\theta}q^*_H
\ge 0.
\]

Assume also that the agent's reservation utility is zero.

Then the participation constraints are
\[
t-\underline{\theta}q\ge 0,
\]
and
\[
t-\overline{\theta}q\ge 0.
\]

Under complete information, the principal offers transfers
\[
t^*(\underline{\theta})
\]
and
\[
t^*(\overline{\theta})
\]
that exactly satisfy the participation constraints.



For the production level
\[
q^*_H,
\]
the transfer is
\[
t^*_H=\overline{\theta}q^*_H
\]
when
\[
\theta=\overline{\theta}.
\]

For the production level
\[
q^*_L,
\]
the transfer is
\[
t^*_L=\underline{\theta}q^*_L
\]
when
\[
\theta=\underline{\theta}.
\]

Whatever his type, the agent accepts the offer and earns zero rent.

The complete-information contracts are therefore
\[
(q^*_L,t^*_L)
\]
for the efficient type and
\[
(q^*_H,t^*_H)
\]
for the inefficient type.

Figure 7 shows the indifference curves of both types.

The agent's utility is
\[
U=t-\theta q.
\]

Thus, higher utility is obtained in the northwest direction in the \((q,t)\)-space.

Since
\[
\overline{\theta}>\underline{\theta},
\]
the indifference curve of the efficient type is flatter than that of the inefficient type.

This is the single-crossing condition, also called the Spence--Mirrlees property.

We can also draw the indifference curves of the principal in the same space.

The principal's payoff is
\[
V=S(q)-t.
\]

Thus, the principal's indifference curves are concave, and the principal's payoff increases in the southeast direction.

The complete-information optimal contract is represented by the pair of points
\[
(A^*,B^*).
\]

For each point, the principal's strictly concave indifference curve is tangent to the zero-rent isoutility curve of the corresponding type.

Therefore,
\[
V^*_L=W^*_L
\]
and
\[
V^*_H=W^*_H
\]
under complete information.

The social optimality conditions are satisfied.

The principal obtains higher profit when dealing with the efficient type, and
\[
q^*_L>q^*_H.
\]

However, the transfer to the efficient type,
\[
t^*_L=\underline{\theta}q^*_L,
\]
is not necessarily greater than the transfer to the inefficient type,
\[
t^*_H=\overline{\theta}q^*_H.
\]

This is because the efficient type has a lower unit cost but produces a larger quantity. The total transfer depends on both marginal cost and output.



\subsection{Incentive Compatibility and Participation}

With private information, there is no way for the principal to distinguish
\[
\overline{\theta}
\]
from
\[
\underline{\theta}.
\]

Hence, both types may select the same contract if the complete-information menu is offered.

That is, both types may prefer the contract intended for the inefficient type to the contract intended for the efficient type.

Therefore, the complete-information menu is generally not incentive compatible.

A menu of contracts
\[
\{(\underline t,\underline q),(\overline t,\overline q)\}
\]
is incentive compatible when:

\begin{itemize}
    \item the efficient type \(\underline{\theta}\) weakly prefers \((\underline t,\underline q)\) to \((\overline t,\overline q)\),
    \item the inefficient type \(\overline{\theta}\) weakly prefers \((\overline t,\overline q)\) to \((\underline t,\underline q)\).
\end{itemize}

The incentive compatibility constraints are therefore
\[
\underline t-\underline{\theta}\underline q
\ge
\overline t-\underline{\theta}\overline q,
\]
and
\[
\overline t-\overline{\theta}\overline q
\ge
\underline t-\overline{\theta}\underline q.
\]

Furthermore, for the menu to be accepted, the participation constraints must be satisfied:
\[
\underline t-\underline{\theta}\underline q
\ge 0,
\]
and
\[
\overline t-\overline{\theta}\overline q
\ge 0.
\]

These four conditions define the set of incentive-feasible allocations achievable through a menu of contracts.

Adding the two incentive compatibility constraints yields
\[
(\overline{\theta}-\underline{\theta})(\underline q-\overline q)\ge 0.
\]

Since
\[
\overline{\theta}-\underline{\theta}>0,
\]
we obtain the monotonicity condition
\[
\underline q\ge \overline q.
\]

Suppose a menu
\[
\{(\underline t,\underline q),(\overline t,\overline q)\}
\]
is offered.

When the efficient type mimics the inefficient type, his payoff is
\[
\overline t-\underline{\theta}\overline q.
\]

This can be rewritten as
\[
\overline t-\underline{\theta}\overline q
=
\overline t-\overline{\theta}\overline q
+
(\overline{\theta}-\underline{\theta})\overline q.
\]

Define
\[
\Delta\theta
=
\overline{\theta}-\underline{\theta}.
\]

Then
\[
\overline t-\underline{\theta}\overline q
=
\overline U+\Delta\theta\,\overline q.
\]

The term
\[
\Delta\theta\,\overline q
\]
is called the information rent.

This information rent is generated by the agent's informational advantage over the principal.

The principal's problem is to determine the least costly way to give up this rent.

We use the notation
\[
\underline U
=
\underline t-\underline{\theta}\underline q,
\]
and
\[
\overline U
=
\overline t-\overline{\theta}\overline q.
\]

\subsection{The Optimal Contract}

The principal's problem is
\[
\max_{\{(\underline t,\underline q),(\overline t,\overline q)\}}
\nu
\left[
S(\underline q)-\underline t
\right]
+
(1-\nu)
\left[
S(\overline q)-\overline t
\right]
\]
subject to the incentive compatibility and participation constraints.

Using
\[
\underline t=\underline U+\underline{\theta}\underline q
\]
and
\[
\overline t=\overline U+\overline{\theta}\overline q,
\]
we can rewrite the problem as
\[
\max_{\{(\underline U,\underline q),(\overline U,\overline q)\}}
\nu
\left[
S(\underline q)-\underline{\theta}\underline q
\right]
+
(1-\nu)
\left[
S(\overline q)-\overline{\theta}\overline q
\right]
-
\left[
\nu\underline U+(1-\nu)\overline U
\right].
\]

The constraints become
\[
\underline U
\ge
\overline U+\Delta\theta\,\overline q,
\]
\[
\overline U
\ge
\underline U-\Delta\theta\,\underline q,
\]
\[
\underline U\ge 0,
\]
and
\[
\overline U\ge 0.
\]

The term
\[
\nu
\left[
S(\underline q)-\underline{\theta}\underline q
\right]
+
(1-\nu)
\left[
S(\overline q)-\overline{\theta}\overline q
\right]
\]
is expected allocative efficiency.

The term
\[
\nu\underline U+(1-\nu)\overline U
\]
is expected information rent.

Thus, the principal maximizes expected social value of trade minus expected information rent.

The principal may accept distortions away from efficiency in order to reduce the agent's information rent.




Note that constraints \((5)\) and \((8)\) together imply
\[
\underline U\ge 0.
\]

Therefore, constraint \((7)\) is irrelevant.

Furthermore, we know that the efficient type wants to mimic the inefficient type, not vice versa.

Therefore, constraint \((6)\) is also irrelevant.

It can be shown that constraints \((5)\) and \((8)\) are binding.

Suppose not, and assume
\[
\overline U=\varepsilon>0.
\]

Then the principal can decrease
\[
\overline U
\]
by \(\varepsilon\) and increase her payoff.

Also, if
\[
\underline U
=
\Delta\theta\,\overline q+\varepsilon,
\qquad
\varepsilon>0,
\]
then the principal can decrease
\[
\underline U
\]
by \(\varepsilon\).

Therefore,
\[
\overline U=0
\]
and
\[
\underline U=\Delta\theta\,\overline q.
\]

Hence, the principal's problem becomes
\[
\max_{\{\underline q,\overline q\}}
\nu
\left[
S(\underline q)-\underline{\theta}\underline q
\right]
+
(1-\nu)
\left[
S(\overline q)-\overline{\theta}\overline q
\right]
-
\nu\Delta\theta\,\overline q.
\]

The first-order conditions are
\[
S'(\underline q^{SB})
=
\underline{\theta},
\]
or equivalently,
\[
\underline q^{SB}=\underline q^*.
\]

For the inefficient type,
\[
(1-\nu)
\left[
S'(\overline q^{SB})-\overline{\theta}
\right]
=
\nu\Delta\theta.
\]

Equivalently,
\[
S'(\overline q^{SB})
=
\overline{\theta}
+
\frac{\nu}{1-\nu}\Delta\theta.
\]

The left-hand side of the condition for \(\overline q^{SB}\) captures the marginal impact of the inefficient agent's output on allocative efficiency.

The right-hand side captures the marginal impact of \(\overline q\) on the information rent paid to the efficient type.

We have
\[
\underline q^{SB}
=
\underline q^*.
\]

Moreover,
\[
\underline q^*>\overline q^*>\overline q^{SB}.
\]

The inequality
\[
\underline q^*>\overline q^*
\]
follows from monotonicity, since
\[
\underline{\theta}<\overline{\theta}.
\]

Also,
\[
S'(\overline q^{SB})-\overline{\theta}>0,
\]
while
\[
S'(\overline q^*)-\overline{\theta}=0.
\]

Since
\[
S''<0,
\]
this implies
\[
\overline q^{SB}<\overline q^*.
\]

\paragraph{Proposition 1.}

Under asymmetric information, the optimal menu of contracts entails the following:

\begin{enumerate}
    \item There is no output distortion for the efficient type:
    \[
    \underline q^{SB}=\underline q^*.
    \]

    \item There is a downward output distortion for the inefficient type:
    \[
    \overline q^{SB}<\overline q^*,
    \]
    with
    \[
    S'(\overline q^{SB})
    =
    \overline{\theta}
    +
    \frac{\nu}{1-\nu}\Delta\theta.
    \]

    \item Only the efficient type receives positive information rent:
    \[
    \underline U^{SB}
    =
    \Delta\theta\,\overline q^{SB},
    \]
    while
    \[
    \overline U^{SB}=0.
    \]

    \item The second-best transfers are
    \[
    \underline t^{SB}
    =
    \underline{\theta}\underline q^*
    +
    \Delta\theta\,\overline q^{SB},
    \]
    and
    \[
    \overline t^{SB}
    =
    \overline{\theta}\overline q^{SB}.
    \]
\end{enumerate}

\paragraph{Exercise 2.}

A businessman wants to contract with a worker.

The worker can be one of two types.

The worker's disutility of producing output \(q\) is either
\[
q^2
\]
or
\[
2q^2.
\]

Therefore, the worker's utility is either
\[
U_G(t,q)=t-q^2,
\]
or
\[
U_B(t,q)=t-2q^2,
\]
depending on his type.

The probability of type \(G\) is
\[
\nu.
\]

Both types have reservation utility equal to zero.

The businessman's utility is
\[
S(q)=kq.
\]

Formulate and solve the businessman's problem under perfect information about the worker's type.

What output levels are demanded, and what wages are paid?

Then formulate and solve the problem when an adverse selection problem is present.

Calculate the optimal contract characteristics.

Compare the cases of symmetric and asymmetric information.




\subsection{An Example: Uniform Nonlinear Pricing}

Uniform nonlinear pricing is also called second-degree price discrimination.

\subsubsection{Setup}

There are two types of consumers.

Type 1 has utility
\[
u_1(x_1)+y_1,
\]
and type 2 has utility
\[
u_2(x_2)+y_2.
\]

Here, \(x\) is the good in question, whose price is set by the monopolist, and \(y\) represents all other goods.

The price of \(y\) is normalized to one.

Assume the single-crossing property:
\[
u_2(0)=u_1(0),
\]
and
\[
u_2'(x)>u_1'(x),
\qquad
\forall x\ge 0.
\]

The monopolist offers a menu of contracts:
\[
(r_i,x_i),
\]
where \(r_i\) is the total payment and \(x_i\) is the quantity received by type \(i\).

\subsubsection{The Monopolist's Problem}

The monopolist solves
\[
\max_{r_1,x_1,r_2,x_2}
\left\{
r_1+r_2-c(x_1+x_2)
\right\}
\]
subject to the participation constraints
\[
u_1(x_1)-r_1\ge 0,
\]
\[
u_2(x_2)-r_2\ge 0,
\]
and the incentive compatibility constraints
\[
u_1(x_1)-r_1
\ge
u_1(x_2)-r_2,
\]
\[
u_2(x_2)-r_2
\ge
u_2(x_1)-r_1.
\]

We can show that constraints \((1)\) and \((4)\) are binding.

First, note that the monopolist wants to make \(r_1\) and \(r_2\) as large as possible.

Only constraints \((1)\) and \((3)\) limit \(r_1\) from becoming arbitrarily large.

Therefore, at least one of \((1)\) and \((3)\) must bind.

Similarly, at least one of \((2)\) and \((4)\) must bind.




Second, note that
\[
u_2(x)>u_1(x),
\qquad
\forall x>0.
\]

Then constraint \((1)\) implies
\[
u_2(x_1)-r_1>0.
\]

By constraint \((4)\),
\[
u_2(x_2)-r_2
\ge
u_2(x_1)-r_1.
\]

Therefore,
\[
u_2(x_2)-r_2>0.
\]

This shows that constraint \((2)\) is not binding.

Consequently, since at least one of \((2)\) and \((4)\) must bind, constraint \((4)\) must be binding.

Third, clearly
\[
r_2>r_1.
\]

Rearranging constraint \((4)\), which is binding, gives
\[
u_2(x_2)-u_2(x_1)
=
r_2-r_1.
\]

This implies
\[
x_2>x_1.
\]

Furthermore, since
\[
u_2'(x)>u_1'(x),
\]
we have
\[
u_2(x_2)-u_2(x_1)
>
u_1(x_2)-u_1(x_1).
\]

Hence,
\[
r_2-r_1
>
u_1(x_2)-u_1(x_1).
\]

Rearranging, we obtain
\[
u_1(x_1)-r_1
>
u_1(x_2)-r_2.
\]

Thus, constraint \((3)\) is not binding.

Therefore, constraint \((1)\) must be binding.

\subsubsection{Solving the Reduced Problem}

By constraint \((1)\),
\[
r_1=u_1(x_1).
\]

By constraints \((1)\) and \((4)\),
\[
r_2
=
u_2(x_2)-u_2(x_1)+u_1(x_1).
\]

Therefore, the original problem reduces to
\[
\max_{x_1,x_2}
\left\{
u_1(x_1)
+
u_2(x_2)
-
u_2(x_1)
+
u_1(x_1)
-
c(x_1+x_2)
\right\}.
\]

Equivalently,
\[
\max_{x_1,x_2}
\left\{
2u_1(x_1)
+
u_2(x_2)
-
u_2(x_1)
-
c(x_1+x_2)
\right\}.
\]

The first-order conditions are
\[
2u_1'(x_1)-u_2'(x_1)-c=0,
\]
and
\[
u_2'(x_2)-c=0.
\]

Thus,
\[
u_1'(x_1)
=
c+\left[u_2'(x_1)-u_1'(x_1)\right].
\]

Since
\[
u_2'(x_1)>u_1'(x_1),
\]
we have
\[
u_1'(x_1)>c.
\]

Hence, type 1 consumes less than the efficient quantity.

For type 2,
\[
u_2'(x_2)=c.
\]

Thus, type 2 consumes the efficient quantity.

The interpretation is that type 2 consumers consume the efficient quantity, while type 1 consumers consume less than the efficient quantity.

Finally, equations \((1)\) and \((4)\) complete the solution:
\[
r_1=u_1(x_1),
\]
and
\[
r_2
=
u_2(x_2)-u_2(x_1)+u_1(x_1).
\]

Moreover,
\[
r_2<u_2(x_2),
\]
so type 2 receives positive information rent.

\paragraph{Exercise 3.}

Let
\[
u_1(x)=4x-\frac{x^2}{4}
\]
and
\[
u_2(x)=8x-\frac{x^2}{2}
\]
be the utility functions of type 1 and type 2 consumers, respectively.

Let marginal cost be
\[
c=2.
\]

There are three type 1 consumers and one type 2 consumer.

The monopolist cannot identify consumer type.

Derive the optimal menu
\[
(r_i,x_i),
\qquad
i=1,2.
\]

\section{Moral Hazard: Hidden Action}

Agents to whom a task has been delegated by a principal may choose actions that affect the value of trade or the agent's performance.

If those actions are unobservable, then they cannot be directly contracted upon.

Unlike adverse selection, uncertainty is now endogenous rather than exogenous.

The probabilities of different states of nature depend explicitly on the agent's effort as well as on pure random variables.

The random output aggregates both the agent's effort and the realization of pure luck.

However, the principal can design a contract only based on the agent's observable performance.

\subsection{The Model}

The agent's effort is
\[
e\in\{0,1\}.
\]

The agent's disutility from effort is
\[
\psi(e),
\]
where
\[
\psi(0)=\psi_0=0,
\]
and
\[
\psi(1)=\psi_1=\psi.
\]

The agent's utility function is
\[
U=u(t)-\psi(e),
\]
where
\[
u'(t)>0,
\qquad
u''(t)<0.
\]




Here,
\[
t
\]
is a transfer from the principal to the agent.

Production is stochastic:
\[
q\in\{\underline q,\overline q\}.
\]

The probability of high output depends on effort:
\[
\Pr(q=\overline q\mid e=0)=\pi_0,
\]
and
\[
\Pr(q=\overline q\mid e=1)=\pi_1.
\]

Assume
\[
\Delta\pi=\pi_1-\pi_0>0.
\]

The transfer \(t\) may or may not depend on realized output \(q\).

The principal is always risk-neutral.

The agent may or may not be risk-neutral.

Effort improves production in the sense of first-order stochastic dominance.

That is,
\[
\Pr(q\le q^*\mid e)
\]
is decreasing in \(e\) for any given production level \(q^*\).

Indeed,
\[
\Pr(q\le \underline q\mid e=1)
=
1-\pi_1
<
1-\pi_0
=
\Pr(q\le \underline q\mid e=0),
\]
and
\[
\Pr(q\le \overline q\mid e=1)
=
1
=
\Pr(q\le \overline q\mid e=0).
\]

\subsection{The Complete Information Optimal Contract}

If the principal wants to induce high effort,
\[
e=1,
\]
then the principal's problem is
\[
\max_{\{\underline t,\overline t\}}
\pi_1(\overline S-\overline t)
+
(1-\pi_1)(\underline S-\underline t)
\]
subject to the participation constraint
\[
\pi_1u(\overline t)
+
(1-\pi_1)u(\underline t)
-
\psi
\ge 0.
\]

Since effort is observable under complete information, no incentive constraint is needed.

The Lagrangian is
\[
\mathcal{L}
=
\pi_1(\overline S-\overline t)
+
(1-\pi_1)(\underline S-\underline t)
+
\lambda
\left[
\pi_1u(\overline t)
+
(1-\pi_1)u(\underline t)
-
\psi
\right].
\]

The first-order conditions are
\[
-\pi_1+\lambda\pi_1u'(\overline t^*)=0,
\]
and
\[
-(1-\pi_1)+\lambda(1-\pi_1)u'(\underline t^*)=0.
\]

Equivalently,
\[
\lambda u'(\overline t^*)=1,
\]
and
\[
\lambda u'(\underline t^*)=1.
\]

Therefore,
\[
u'(\overline t^*)=u'(\underline t^*).
\]

Since \(u\) is strictly concave, \(u'\) is strictly decreasing. Hence,
\[
\overline t^*=\underline t^*.
\]

Thus, under complete information, the agent is fully insured against output risk.

The participation constraint binds:
\[
\pi_1u(\overline t^*)
+
(1-\pi_1)u(\underline t^*)
-
\psi
=
0.
\]

Since
\[
\overline t^*=\underline t^*,
\]
we obtain
\[
u(t^*)=\psi,
\]
or
\[
t^*=u^{-1}(\psi).
\]



From the first-order conditions, we have
\[
\lambda
=
\frac{1}{u'(\overline t^*)}
=
\frac{1}{u'(\underline t^*)}
>0.
\]

Therefore,
\[
\overline t^*=\underline t^*=t^*.
\]

Hence, the agent obtains full insurance.

Define the inverse utility function
\[
h=u^{-1},
\]
where
\[
h'>0,
\qquad
h''>0.
\]

Since the participation constraint is binding, we obtain the value of the transfer that is just enough to cover the disutility of effort:
\[
t^*=h(\psi).
\]

This is also the first-best expected payment to the agent:
\[
C^{FB}=h(\psi).
\]

Which level of effort is preferred by the principal?

When
\[
e=1,
\]
the principal's expected payoff is
\[
V_1
=
\pi_1\overline S+(1-\pi_1)\underline S-h(\psi).
\]

When
\[
e=0,
\]
the principal's expected payoff is
\[
V_0
=
\pi_0\overline S+(1-\pi_0)\underline S.
\]

Inducing effort is optimal from the principal's point of view when
\[
V_1\ge V_0.
\]

Equivalently,
\[
\Delta\pi\Delta S\ge h(\psi),
\]
where
\[
\Delta\pi=\pi_1-\pi_0
\]
and
\[
\Delta S=\overline S-\underline S.
\]

Here,
\[
\Delta\pi\Delta S
\]
is the principal's expected gain from effort, while
\[
h(\psi)
\]
is the first-best cost of inducing effort.

\subsection{Risk-Neutral Agents}

Suppose the agent is risk-neutral.

Then
\[
u(t)=t
\]
for all \(t\), and therefore
\[
h(u)=u.
\]

The principal who wants to induce effort solves
\[
\max_{\{\overline t,\underline t\}}
\pi_1(\overline S-\overline t)
+
(1-\pi_1)(\underline S-\underline t)
\]
subject to
\[
\pi_1\overline t+(1-\pi_1)\underline t-\psi
\ge
\pi_0\overline t+(1-\pi_0)\underline t,
\]
and
\[
\pi_1\overline t+(1-\pi_1)\underline t-\psi
\ge 0.
\]

Now the incentive compatibility constraint is also needed.

With risk neutrality, the principal can choose incentive-compatible transfers
\[
\overline t
\quad
\text{and}
\quad
\underline t
\]
such that both constraints bind.

Solving the two binding constraints gives
\[
\overline t^*
=
\frac{1-\pi_0}{\Delta\pi}\psi,
\]
and
\[
\underline t^*
=
-\frac{\pi_0}{\Delta\pi}\psi.
\]

The expected payment by the principal is
\[
\pi_1\overline t^*
+
(1-\pi_1)\underline t^*
=
\psi.
\]

Thus, the expected payment is exactly equal to the disutility of effort.

By increasing effort from
\[
e=0
\]
to
\[
e=1,
\]
the agent receives the high-output transfer
\[
\overline t^*
\]
more often than the low-output transfer
\[
\underline t^*.
\]

Since
\[
\Delta\pi>0,
\]
the agent's expected gain from exerting effort is
\[
\Delta\pi(\overline t^*-\underline t^*)
=
\psi.
\]

This exactly compensates the agent for the additional disutility of effort.

Therefore, there is no loss of efficiency.

The principal can induce high effort without reducing social welfare.

\paragraph{Theorem 3.}

Moral hazard is not an issue with a risk-neutral agent despite the nonobservability of effort.

The first-best level of effort is still implemented.



\paragraph{Exercise 4.}

The agent is risk-neutral but has no wealth.

The principal cannot impose liability on the agent even if the outcome is very bad.

Thus, payments to the agent cannot be negative:
\[
\overline t\ge 0,
\qquad
\underline t\ge 0.
\]

Derive the expected payment that must be allowed to the agent in order to induce high effort under these limited-liability constraints.

How do these constraints affect the efficiency of the contract?

\subsection{Risk Aversion and Optimal Contract}

When the agent is risk-averse, the principal's problem is
\[
\max_{\{\overline t,\underline t\}}
\pi_1(\overline S-\overline t)
+
(1-\pi_1)(\underline S-\underline t)
\]
subject to the incentive compatibility constraint
\[
\pi_1u(\overline t)
+
(1-\pi_1)u(\underline t)
-
\psi
\ge
\pi_0u(\overline t)
+
(1-\pi_0)u(\underline t),
\]
and the participation constraint
\[
\pi_1u(\overline t)
+
(1-\pi_1)u(\underline t)
-
\psi
\ge 0.
\]

Since the concave function \(u(\cdot)\) appears on both sides of the incentive compatibility constraint, it is not immediate that the problem is concave.

Define
\[
\overline u=u(\overline t),
\qquad
\underline u=u(\underline t).
\]

Equivalently, let
\[
\overline t=h(\overline u),
\qquad
\underline t=h(\underline u),
\]
where
\[
h=u^{-1}.
\]

Then the two constraints become
\[
\pi_1\overline u
+
(1-\pi_1)\underline u
-
\psi
\ge
\pi_0\overline u
+
(1-\pi_0)\underline u,
\]
and
\[
\pi_1\overline u
+
(1-\pi_1)\underline u
-
\psi
\ge 0.
\]

The transformed problem is
\[
\max_{\{\overline u,\underline u\}}
\pi_1\left[\overline S-h(\overline u)\right]
+
(1-\pi_1)\left[\underline S-h(\underline u)\right]
\]
subject to the two constraints above.

The objective function is strictly concave in
\[
(\overline u,\underline u),
\]
because
\[
h''>0.
\]

Let
\[
\lambda
\]
and
\[
\mu
\]
be the nonnegative multipliers associated with the incentive compatibility and participation constraints, respectively.

The first-order conditions are
\[
-\pi_1h'(\overline u^{SB})
+
\lambda\Delta\pi
+
\mu\pi_1
=
0,
\]
and
\[
-(1-\pi_1)h'(\underline u^{SB})
-
\lambda\Delta\pi
+
\mu(1-\pi_1)
=
0.
\]

Using
\[
h'(u(t))=\frac{1}{u'(t)},
\]
these can be written as
\[
-\frac{\pi_1}{u'(\overline t^{SB})}
+
\lambda\Delta\pi
+
\mu\pi_1
=
0,
\]
and
\[
-\frac{1-\pi_1}{u'(\underline t^{SB})}
-
\lambda\Delta\pi
+
\mu(1-\pi_1)
=
0.
\]

Rearranging terms gives
\[
\frac{1}{u'(\overline t^{SB})}
=
\mu
+
\lambda\frac{\Delta\pi}{\pi_1},
\]
and
\[
\frac{1}{u'(\underline t^{SB})}
=
\mu
-
\lambda\frac{\Delta\pi}{1-\pi_1}.
\]

The four variables
\[
\overline t^{SB},
\quad
\underline t^{SB},
\quad
\lambda,
\quad
\mu
\]
are obtained simultaneously from the two constraints and the two first-order conditions.

Multiplying
\[
\frac{1}{u'(\overline t^{SB})}
=
\mu
+
\lambda\frac{\Delta\pi}{\pi_1}
\]
by \(\pi_1\), and
\[
\frac{1}{u'(\underline t^{SB})}
=
\mu
-
\lambda\frac{\Delta\pi}{1-\pi_1}
\]
by \(1-\pi_1\), then adding, gives
\[
\mu
=
\frac{\pi_1}{u'(\overline t^{SB})}
+
\frac{1-\pi_1}{u'(\underline t^{SB})}
>0.
\]

Hence, the participation constraint is binding.



Similarly, from the first-order conditions we obtain
\[
\lambda
=
\frac{\pi_1(1-\pi_1)}{\Delta\pi}
\left[
\frac{1}{u'(\overline t^{SB})}
-
\frac{1}{u'(\underline t^{SB})}
\right].
\]

The incentive compatibility constraint implies
\[
\overline u^{SB}-\underline u^{SB}
\ge
\frac{\psi}{\Delta\pi}>0.
\]

Therefore,
\[
\overline t^{SB}>\underline t^{SB}.
\]

Since
\[
u''<0,
\]
we have
\[
u'(\overline t^{SB})<u'(\underline t^{SB}).
\]

Thus,
\[
\frac{1}{u'(\overline t^{SB})}
>
\frac{1}{u'(\underline t^{SB})}.
\]

Hence,
\[
\lambda>0.
\]

Therefore, the incentive compatibility constraint is also binding.

Because both constraints are binding, we can obtain the values of
\[
u(\overline t^{SB})
\]
and
\[
u(\underline t^{SB})
\]
by solving
\[
\pi_1\overline u^{SB}
+
(1-\pi_1)\underline u^{SB}
-
\psi
=
0,
\]
and
\[
\Delta\pi
\left(
\overline u^{SB}-\underline u^{SB}
\right)
=
\psi.
\]

Solving these two equations gives
\[
\overline u^{SB}
=
\psi
+
(1-\pi_1)\frac{\psi}{\Delta\pi},
\]
and
\[
\underline u^{SB}
=
\psi
-
\pi_1\frac{\psi}{\Delta\pi}.
\]

Therefore, second-best transfers are
\[
\overline t^{SB}
=
h\left(
\psi
+
(1-\pi_1)\frac{\psi}{\Delta\pi}
\right),
\]
and
\[
\underline t^{SB}
=
h\left(
\psi
-
\pi_1\frac{\psi}{\Delta\pi}
\right).
\]

Since
\[
\overline t^{SB}\ne \underline t^{SB},
\]
the risk-averse agent does not receive full insurance.

Inducing effort requires the principal to let the agent bear some risk.

\paragraph{Theorem 4.}

When the agent is strictly risk-averse, the optimal contract that induces effort makes both the agent's participation constraint and incentive compatibility constraint binding.

This contract does not provide full insurance.

The second-best transfers are
\[
\overline t^{SB}
=
h\left(
\psi
+
(1-\pi_1)\frac{\psi}{\Delta\pi}
\right),
\]
and
\[
\underline t^{SB}
=
h\left(
\psi
-
\pi_1\frac{\psi}{\Delta\pi}
\right).
\]

Note that
\[
\overline t^{SB}>h(\psi)
\]
and
\[
\underline t^{SB}<h(\psi).
\]

That is, the agent receives more than the complete-information transfer when high output is realized, and less than the complete-information transfer when low output is realized.

The second-best cost of inducing effort under moral hazard is
\[
C^{SB}
=
\pi_1\overline t^{SB}
+
(1-\pi_1)\underline t^{SB}.
\]

Therefore,
\[
C^{SB}
=
\pi_1
h\left(
\psi
+
(1-\pi_1)\frac{\psi}{\Delta\pi}
\right)
+
(1-\pi_1)
h\left(
\psi
-
\pi_1\frac{\psi}{\Delta\pi}
\right).
\]

The benefit of inducing effort is still
\[
\Delta\pi\Delta S.
\]

Thus, high effort \(e^*=1\) is optimal whenever
\[
\Delta\pi\Delta S
\ge
C^{SB}.
\]

By convexity of \(h\) and Jensen's inequality,
\[
C^{SB}\ge h(\psi)=C^{FB}.
\]

Therefore, inducing high effort occurs less often under moral hazard than when effort is observable.

\subsubsection{Example with Quadratic Inverse Utility}

Assume
\[
h(u)=u+\frac{r}{2}u^2,
\qquad
r>0.
\]

Let \(u^{SB}\) denote the random utility level that the agent receives across states.

Then
\[
C^{SB}
=
E[h(u^{SB})].
\]

Since
\[
h(u)=u+\frac{r}{2}u^2,
\]
we have
\[
C^{SB}
=
E[u^{SB}]
+
\frac{r}{2}E[(u^{SB})^2].
\]

Equivalently,
\[
C^{SB}
=
E[u^{SB}]
+
\frac{r}{2}
\left[
(E[u^{SB}])^2
+
\operatorname{Var}(u^{SB})
\right].
\]

Since the participation constraint is binding,
\[
E[u^{SB}]
=
\pi_1\overline u^{SB}
+
(1-\pi_1)\underline u^{SB}
=
\psi.
\]

From Theorem 4,
\[
\overline u^{SB}
=
\psi
+
(1-\pi_1)\frac{\psi}{\Delta\pi},
\]
and
\[
\underline u^{SB}
=
\psi
-
\pi_1\frac{\psi}{\Delta\pi}.
\]

Hence,
\[
\operatorname{Var}(u^{SB})
=
\pi_1(1-\pi_1)
\frac{\psi^2}{(\Delta\pi)^2}.
\]

Therefore,
\[
C^{SB}
=
\psi
+
\frac{r}{2}\psi^2
+
\frac{r}{2}
\pi_1(1-\pi_1)
\frac{\psi^2}{(\Delta\pi)^2}.
\]

The first-best cost of implementing effort is
\[
C^{FB}
=
h(\psi)
=
\psi+\frac{r}{2}\psi^2.
\]

Hence, the agency cost is
\[
AC
=
C^{SB}-C^{FB}.
\]

Thus,
\[
AC
=
\frac{r}{2}
\pi_1(1-\pi_1)
\frac{\psi^2}{(\Delta\pi)^2}.
\]

This agency cost is the principal's loss between the first-best and second-best expected profit when implementing positive effort.




The agency cost
\[
AC
=
\frac{r}{2}
\pi_1(1-\pi_1)
\frac{\psi^2}{(\Delta\pi)^2}
\]
increases in:

\begin{itemize}
    \item \(r\), which measures the agent's risk aversion;
    \item \(\psi\), which is the cost of effort;
    \item
    \[
    \eta
    =
    \frac{\pi_1(1-\pi_1)}{(\Delta\pi)^2},
    \]
    which measures the severity of the informational problem faced by the principal.
\end{itemize}

\paragraph{Exercise 5.}

Assume that uncertainty is represented by three states of nature.

The agent can choose between two effort levels:
\[
e=6
\]
and
\[
e=4.
\]

The outcome matrix is
\[
\begin{array}{c|ccc}
 & E_1 & E_2 & E_3 \\
\hline
e=6 & 60000 & 60000 & 30000 \\
e=4 & 30000 & 60000 & 30000
\end{array}
\]

The probability of each state is
\[
\frac{1}{3}.
\]

The objective function of the principal is
\[
B(q,t)=q-t,
\]
and the agent's utility function is
\[
U(t,e)=\sqrt{t}-e^2.
\]

Answer the following questions:

\begin{enumerate}
    \item What are the effort and wages under symmetric information?
    
    \item What happens under asymmetric information?
    
    \item What payoff scheme induces effort level \(e=4\)?
    
    \item What payoff scheme induces effort level \(e=6\)?
    
    \item Which effort level does the principal prefer?
\end{enumerate}

\subsection{A Continuum of Performances}

The level of performance
\[
q
\]
is drawn from a continuous distribution with cumulative distribution function
\[
F(\cdot\mid e)
\]
on the support
\[
[\underline q,\overline q].
\]

This distribution is conditional on effort level
\[
e\in\{0,1\}.
\]

Let
\[
f(\cdot\mid e)
\]
be the conditional density function.

A contract
\[
t(q)
\]
that induces positive effort must satisfy the incentive compatibility constraint
\[
\int_{\underline q}^{\overline q}
u(t(q))f(q\mid 1)\,dq
-
\psi
\ge
\int_{\underline q}^{\overline q}
u(t(q))f(q\mid 0)\,dq,
\]
and the participation constraint
\[
\int_{\underline q}^{\overline q}
u(t(q))f(q\mid 1)\,dq
-
\psi
\ge 0.
\]

The principal's problem is
\[
\max_{\{t(q)\}}
\int_{\underline q}^{\overline q}
\left[
S(q)-t(q)
\right]
f(q\mid 1)\,dq
\]
subject to the incentive compatibility and participation constraints above.

Let
\[
\lambda
\]
and
\[
\mu
\]
be the multipliers on the incentive compatibility and participation constraints, respectively.

The Lagrangian density can be written as
\[
\mathcal{L}(q,t)
=
[S(q)-t(q)]f(q\mid 1)
+
\lambda
\left[
u(t(q))(f(q\mid 1)-f(q\mid 0))
-\psi
\right]
+
\mu
\left[
u(t(q))f(q\mid 1)
-\psi
\right].
\]

Optimizing pointwise with respect to \(t(q)\) yields
\[
-f(q\mid 1)
+
\lambda u'(t^{SB}(q))
\left[
f(q\mid 1)-f(q\mid 0)
\right]
+
\mu u'(t^{SB}(q))f(q\mid 1)
=
0.
\]

Dividing by
\[
f(q\mid 1)>0,
\]
we obtain
\[
\frac{1}{u'(t^{SB}(q))}
=
\mu
+
\lambda
\left[
1-
\frac{f(q\mid 0)}{f(q\mid 1)}
\right].
\]

Equivalently,
\[
u'(t^{SB}(q))
=
\frac{1}{
\mu
+
\lambda
\left[
1-
\dfrac{f(q\mid 0)}{f(q\mid 1)}
\right]
}.
\]

This condition shows that the optimal transfer depends on the likelihood ratio
\[
\frac{f(q\mid 0)}{f(q\mid 1)}.
\]

Outcomes that are relatively more likely under high effort than under low effort should receive higher rewards.



Multiplying the first-order condition by \(f(q\mid 1)\) and taking expectations gives
\[
\mu
=
E_1\left[
\frac{1}{u'(t^{SB}(q))}
\right]
>0,
\]
where \(E_1(\cdot)\) denotes the expectation operator with respect to the probability distribution of output induced by high effort,
\[
e^{SB}=1.
\]

Finally, using this expression for \(\mu\), substituting it into the first-order condition, and multiplying by
\[
f(q\mid 1)u(t^{SB}(q)),
\]
we obtain
\[
\lambda
\bigl[
f(q\mid 1)-f(q\mid 0)
\bigr]
u(t^{SB}(q))
=
f(q\mid 1)u(t^{SB}(q))
\left[
\frac{1}{u'(t^{SB}(q))}
-
E_1\left(
\frac{1}{u'(t^{SB}(q))}
\right)
\right].
\]

Integrating over
\[
[\underline q,\overline q]
\]
and using the complementary slackness condition,
\[
\lambda
\left[
\int_{\underline q}^{\overline q}
\bigl(f(q\mid 1)-f(q\mid 0)\bigr)
u(t^{SB}(q))\,dq
-
\psi
\right]
=
0,
\]
we obtain
\[
\lambda\psi
=
\operatorname{Cov}_1
\left(
u(t^{SB}(q)),
\frac{1}{u'(t^{SB}(q))}
\right)
\ge 0.
\]

Hence,
\[
\lambda\ge 0.
\]

Moreover, \(\lambda=0\) only if \(t^{SB}(q)\) is constant.

However, a constant wage
\[
t^{SB}(q)
\]
cannot satisfy the incentive compatibility constraint when positive effort must be induced.

Therefore,
\[
\lambda>0.
\]

The left-hand side of the first-order condition is increasing in
\[
t^{SB}(q)
\]
because \(u(\cdot)\) is concave.

If
\[
\frac{d}{dq}
\left[
1-
\frac{f(q\mid 0)}{f(q\mid 1)}
\right]
>0,
\]
then
\[
t^{SB}(q)
\]
is monotonically increasing in \(q\).

Equivalently, this condition can be written as
\[
\frac{d}{dq}
\left[
\frac{f(q\mid 1)}{f(q\mid 0)}
\right]
\ge 0.
\]

This property is called the \textbf{monotone likelihood ratio property}.

It is important to note that the monotone likelihood ratio property is not necessarily satisfied even if the first-order stochastic dominance condition holds:
\[
F(q\mid 1)\le F(q\mid 0),
\qquad
\forall q\in[\underline q,\overline q].
\]






